% Acronyms
\newacronym{garch}{GARCH}{Generalized Autoregressive Conditional Heteroskedasticity}

\newacronym{arch}{ARCH}{Autoregressive Conditional Heteroskedasticity}

\newacronym{gjr-garch}{GJR-GARCH}{Glosten-Jagannathan-Runkle GARCH}

\newacronym{egarch}{EGARCH}{Exponential Generalized Autoregressive Conditional Heteroskedasticity}

\newacronym{adf}{ADF}{Augmented Dickey-Fuller}

\newacronym{var}{VaR}{Value at Risk}

\newacronym{nic}{NIC}{News Impact Curve}

\newacronym{ar}{AR}{Autoregressive}

\newacronym{arma}{ARMA}{Autoregressive Moving Average}

\newacronym{iid}{i.i.d.}{Indipendenti e Identicamente Distribuite}

% Glossary entries
\newglossaryentry{scale-free} {
    name=scale-free,
    text=scale-free,
    sort=scale-free,
    description={in ambito finanziario ed econometrico, una misura si definisce \emph{scale-free} (letteralmente ''priva di scala'') quando il suo valore è indipendente dalla scala di grandezza del capitale investito e dal prezzo nominale dell'asset. I rendimenti godono di questa proprietà poiché esprimono una variazione relativa (o percentuale), permettendo di confrontare direttamente la performance di strumenti finanziari o portafogli che presentano valori assoluti profondamente diversi tra loro}
}

\newglossaryentry{omoschedasticita} {
    name=omoschedasticità,
    text=omoschedasticità,
    sort=omoschedasticita,
    description={proprietà di un modello statistico in cui la varianza dell'errore si mantiene costante nel tempo, ipotesi spesso violata nelle serie storiche finanziarie}
}

\newglossaryentry{stazionarieta} {
    name=stazionarietà,
    text=stazionarietà,
    sort=stazionarieta,
    description={condizione di un processo stocastico le cui proprietà statistiche principali (media, varianza e struttura di autocorrelazione) non variano nel tempo}
}

\newglossaryentry{proxy} {
    name=proxy,
    text=proxy,
    sort=proxy,
    description={variabile osservabile e misurabile che viene utilizzata in un modello statistico in sostituzione di una variabile latente o non direttamente quantificabile}
}

\newglossaryentry{variabile-latente} {
    name=variabile latente,
    text=variabile-latente,
    sort=variabile-latente,
    description={variabile che non può essere osservata o misurata direttamente (come la volatilità), ma la cui dinamica viene dedotta tramite l'osservazione di altre variabili correlate}
}